
\documentclass[11pt,a4paper]{article}

\usepackage[utf8]{inputenc}
\usepackage[T1]{fontenc}
\usepackage[french]{babel}
\usepackage{lmodern}
\usepackage[margin=2cm]{geometry}
\usepackage{enumitem}
\usepackage[hidelinks]{hyperref}

% Mise en page 
\setlist[description]{leftmargin=2.8cm, labelwidth=2.5cm, labelsep=0.3cm,
  align=left, font=\normalfont\bfseries, itemsep=0.5em}
\setlist[itemize]{topsep=0.2em, itemsep=0em}
\setlength{\parindent}{0pt}
\pagestyle{empty}

\title{CV}
\author{Victor De Bono}
\date{2026}

\begin{document}

\begin{center}
  {\LARGE\textbf{Victor De Bono}}\\[0.6em]
\end{center}

% ------------------------------------------------------------------
\textbf{Maître de conférences en sciences et techniques pour l'architecture (STA)}, ENSAP Lille\\
\textbf{Chercheur}, LACTH : Laboratoire d'architecture, conception, territoire, histoire, matérialité\\
\textbf{Chercheur associé}, UMR MAP-ARIA : Modèles et simulations pour l'architecture et le patrimoine


% ------------------------------------------------------------------

% ------------------------------------------------------------------
\section*{Recherche}

\subsection*{Articles de revues scientifiques}

\begin{description}
  \item[2026] \textbf{Design, fabrication and testing of a 3D printed concrete shell with continuous fiber-reinforcement}\\
  \textbf{V. De Bono}, R. Mesnil, N. Ducoulombier, J.-F. Caron.\\
  Preprint, \href{https://papers.ssrn.com/sol3/papers.cfm?abstract_id=6213618}{SSRN 6213618}

  \item[2026] \textbf{Fast penetration testing of printable concretes with a portable device: robustness and calibration}\\
  L. Caneda-Martínez, M. Hassan, L. Demont, E. Keita, P. Belin, \textbf{V. De Bono}, R. Mesnil, J.-F. Caron, B. González-Fonteboa, N. Roussel.\\
  \emph{Cement and Concrete Research}, \href{https://doi.org/10.1016/j.cemconres.2026.108141}{doi:10.1016/j.cemconres.2026.108141}

  \item[2025] \textbf{Mechanical properties of 3D printed concrete: a RILEM 304-ADC interlaboratory study -- compressive strength and modulus of elasticity}\\
  V. Mechtcherine, S. Muthukrishnan, A. Robens-Radermacher, R. Wolfs, J. Versteege, C. Menna, O. Ozturk, N. Ozyurt, J. Roupec, C. Richter, J. Jungwirth, L. Miranda, R. Ammann, J.-F. Caron, \textbf{V. De Bono}, et al.\\
  \emph{Materials and Structures}, \href{https://doi.org/10.1617/s11527-025-02688-9}{doi:10.1617/s11527-025-02688-9}

  \item[2025] \textbf{Mechanical properties of 3D printed concrete: a RILEM TC 304-ADC interlaboratory study -- flexural and tensile strength}\\
  V. Mechtcherine, S. Muthukrishnan, A. Robens-Radermacher, R. Wolfs, J. Versteege, C. Menna, O. Ozturk, N. Ozyurt, J. Roupec, C. Richter, J. Jungwirth, L. Miranda, R. Ammann, J.-F. Caron, \textbf{V. De Bono}, et al.\\
  \emph{Materials and Structures}, \href{https://doi.org/10.1617/s11527-025-02687-w}{doi:10.1617/s11527-025-02687-w}

  \item[2024] \textbf{Snapshot on 3D printing with alternative binders and materials: Earth, geopolymers, gypsum and low carbon concrete}\\
  A. Perrot, Y. Jacquet, J.-F. Caron, R. Mesnil, N. Ducoulombier, \textbf{V. De Bono}, J. Sanjayan, S. Ramakrishnan, H. Kloft, J. Gosslar, S. Muthukrishnan, V. Mechtcherine, T. Wangler, J. L. Provis, K. Dörfler, E. Krakovska, N. Roussel, E. Keita.\\
  \emph{Cement and Concrete Research}, \href{https://doi.org/10.1016/j.cemconres.2024.107651}{doi:10.1016/j.cemconres.2024.107651}

  \item[2024] \textbf{Methodology for formulating low-carbon printable mortar through particles packing optimization}\\
  \textbf{V. De Bono}, N. Ducoulombier, R. Mesnil, J.-F. Caron.\\
  \emph{Cement and Concrete Research}, \href{https://doi.org/10.1016/j.cemconres.2023.107403}{doi:10.1016/j.cemconres.2023.107403}
\end{description}

\subsection*{Conférences}

\begin{description}
  \item[2026] \textbf{Conception et fabrication d'une coque imprimée en 3D renforcée par des fibres continues}, SCAN'26, ENSA Montpellier\\
  \textbf{V. De Bono}, R. Mesnil, N. Ducoulombier, J.-F. Caron.\\
  En cours de publication.

  \item[2024] \textbf{Design, structural optimization and fabrication of concrete shell through fiber-reinforced 3D printing}, IASS Symposium, ETH Zurich\\
  \textbf{V. De Bono}, R. Mesnil, N. Ducoulombier, J.-F. Caron.\\
  \href{https://www.ingentaconnect.com/contentone/iass/piass/2024/00002024/00000004/art00009}{Proceedings of the IASS Annual Symposium 2024}

  \item[2024] \textbf{Toward 3D printable low carbon mortar. Method and application}, RILEM, Milan\\
  \textbf{V. De Bono}, S. Loulha, N. Ducoulombier, R. Mesnil, J.-F. Caron.\\
  \href{https://link.springer.com/chapter/10.1007/978-3-031-70281-5_25}{doi:10.1007/978-3-031-70281-5\_25}

  \item[2023] \textbf{Formulation de béton pour l'impression 3D}, AFJC, Paris\\
  \textbf{V. De Bono}, N. Ducoulombier, R. Mesnil, J.-F. Caron.\\
  \href{https://doi.org/10.26168/ajce.41.1.4}{doi:10.26168/ajce.41.1.4}

  \item[2023] \textbf{3D printing of continuous-fibers cementitious composites}, Strategies for Reinforcing Additively Manufactured Constructions (en ligne)\\
  \textbf{V. De Bono}, L. Demont, N. Ducoulombier, R. Mesnil, J.-F. Caron.\\
  \href{https://doi.org/10.52825/ocp.v3i.193}{doi:10.52825/ocp.v3i.193}
\end{description}

\subsection*{Ouvrage}

\begin{description}
  \item[2025] \textbf{Impression 3D béton : état de l'art et applications}, A. Perrot et Y. Jacquet (dir.), ISTE Éditions, ISBN 9781789482003\\
  Co-écriture du chapitre 3, \emph{Caractérisation des matériaux imprimés frais et durcis, du laboratoire à la pratique} : \textbf{V. De Bono}, N. Ducoulombier, A. Perrot.
\end{description}

\subsection*{Thèse}

\begin{description}
  \item[2025] \textbf{Étude des procédés d'impression 3D de bétons fibrés pour la construction. Formulation, renforcement et application}\\
  Sous la direction de J.-F. Caron, co-encadrée par N. Ducoulombier et R. Mesnil.\\
  \href{https://theses.hal.science/tel-05329680/}{theses.hal.science/tel-05329680}
\end{description}


\section*{Enseignement}

\begin{description}
  \item[2026 -- auj.] \textbf{Maître de conférences titulaire STA-CIMA}, ENSAP Lille
  \begin{itemize}
    \item TD Franchissement, licence 1 - 36 h TD, 25 étudiants.
    \item CM et TD Thermique du bâtiment, licence 2 -- 22 h CM, 24 h TD
    \item CM et TD Système constructif, licence 2 - 22 h CM, 24h TD
    \item TD Architecture bioclimatique, licence 3 - 12 h TD
    \item TD Enveloppes, licence 3 - 24 h TD
    \item CM Analyse de cycle de vie, master 1 - 22 h CM
    \item Exploratoire Fabrication additive, master 1 - 48 h TD
  \end{itemize}

  \item[2025 -- 2026] \textbf{Maître de conférences associé STA}, ENSA Bretagne
  \begin{itemize}
    \item TD Cultures matérielles et constructives, licence 1 - 24 h TD
    \item TD Modélisation BIM, Revit, licences 2 et 3 - 96 h TD
    \item CM Thermique du bâtiment, licence 2 - 24 h TD
    \item Atelier Architecte-ingénieur, master 1 - 96 h TD
    \item TD Modélisation BIM et paramétrique, Revit/Dynamo, master 1 - 48 h TD
    \item Approfondissement \emph{Hybridations : artefact, calcul, structure}, master 1 - 12 h TD
  \end{itemize}

  \item[2019 -- 2025] \textbf{Enseignant vacataire}, ENSA Lyon
  \begin{itemize}
    \item Intensif numérique, modélisation paramétrique (Rhinoceros, Grasshopper), licences 2 et 3.
    \item CM Continuum numérique et cartographie du numérique, licence 2 et master FPC.
  \end{itemize}

  \item[2023 -- 2024] \textbf{Enseignant vacataire}, ENSA Lyon
  \begin{itemize}
    \item Accompagnement de projet Archifolies, master 1.
  \end{itemize}

  \item[2022 -- 2025] \textbf{Enseignant vacataire}, École nationale des ponts et chaussées
  \begin{itemize}
    \item TP Connaissance des matériaux, 1\textsuperscript{re}, 2\textsuperscript{e} et 3\textsuperscript{e} années.
    \item Workshop impression 3D béton et argile pour la construction, mastère spécialisé \emph{Design by Data}.
  \end{itemize}

  \item[2018 -- 2021] \textbf{Enseignant vacataire}, ENTPE
  \begin{itemize}
    \item Outils de modélisation BIM, 3\textsuperscript{e} année.
  \end{itemize}

  \item[2022 -- 2024] \textbf{Intervenant}, ÉTS Montréal et Faculté de l'aménagement, Université de Montréal
  \begin{itemize}
    \item CM Fabrication additive et construction : historique, technologies et applications, 2\textsuperscript{e} année.
  \end{itemize}

  \item[2021 -- 2022] \textbf{Workshop}, École Camondo
  \begin{itemize}
    \item Outils de conception paramétrique et impression 3D pour l'architecture intérieure, 1\textsuperscript{re} à 3\textsuperscript{e} années.
  \end{itemize}

  \item[2018 -- 2020] \textbf{Enseignant vacataire}, Institut d'urbanisme de Lyon, Université Lyon 2
  \begin{itemize}
    \item Outils de représentation pour l'urbanisme (suite Adobe), licence 2.
  \end{itemize}
\end{description}

% ------------------------------------------------------------------
\section*{Projets (XtreeE)}

\begin{description}
  \item[2025] \textbf{Portiques structurels}\\
  Conception et fabrication d'éléments de structure poteaux-poutres fonctionnant en portique (encastrement). Mise au point du système constructif, intégration des CET, phases ESQ à EXE.\\
  Maîtrise d'œuvre : XtreeE, Bouygues Construction, Scale One. Maîtrise d'ouvrage : confidentiel.

  \item[2020 -- 2025] \textbf{Passerelle Cabrié}\\
  Conception et fabrication d'une passerelle piétonne à Clermont-l'Hérault en conception-réalisation, phases ESQ à EXE.\\
  Maîtrise d'œuvre : XD Architecture, XtreeE, Sedoa. Maîtrise d'ouvrage : Département de l'Hérault. Construction : Freyssinet, Spie Batignolles, XtreeE.

  \item[2021 -- 2024] \textbf{Passerelle Plaine Commune}\\
  Conception d'une passerelle pour les Jeux olympiques de Paris, projet arrêté en phase 2 de R\&D. Études ACV et certification avec le CSTB, phases ESQ à APD.\\
  Maîtrise d'œuvre : Lavigne \& Cheron, XtreeE, Sedoa. Maîtrise d'ouvrage : Solideo. Construction : Freyssinet, XtreeE.

  \item[2021] \textbf{Coffrages perdus}\\
  Conception et fabrication de coffrages perdus pour faciliter la mise en œuvre : coffrages de nœuds de poutres et boîtes de réservation.\\
  Conception : XtreeE. Client : Spie Batignolles.

  \item[2021 -- 2023] \textbf{Récifs artificiels}\\
  Conception et fabrication de cinq typologies de récifs artificiels.\\
  Conception : XtreeE. Clients : Seaboost et Setec In Vivo.

  \item[2019 -- 2022] \textbf{Viliaprint, maisons à Reims}\\
  Cinq maisons imprimées en 3D béton, certifiées par une appréciation technique d'expérimentation (ATEx), phases APS à EXE.\\
  Maîtrise d'ouvrage : Plurial Novilia. Maîtrise d'œuvre : Emmanuel Coste, XtreeE, NECS. Réalisation : XtreeE et Demathieu Bard.

  \item[2019 -- 2021] \textbf{Archipel, aménagement paysager}\\
  Aménagement de trois cours intérieures du siège social de Vinci Construction, phases ESQ à EXE.\\
  Maîtrise d'ouvrage : Vinci Construction France. Maîtrise d'œuvre : Viguier. Réalisation : XtreeE et VCF.

  \item[2019 -- 2021] \textbf{Mât de télécommunication}\\
  Mât support de télécommunication végétalisé en Guadeloupe, phases ESQ à EXE.\\
  Maîtrise d'ouvrage : Art \& Fact Innovation. Conception : XtreeE et Lamoureux \& Ricciotti Ingénierie. Végétalisation : Jardins de Babylone. Fabrication : XtreeE. Mise en œuvre : Freyssinet.
\end{description}

\section*{Expérience professionnelle}

\begin{description}
  \item[2019 -- 2025] \textbf{Chef de projet}, XtreeE
  \begin{itemize}
    \item Gestion de projets d'impression 3D grande échelle pour la construction, de la conception à la livraison.
    \item Conception et modélisation paramétrique appliquées aux éléments architecturaux et structurels, phases APS, APD.
    \item Réalisation de carnets de détails, systèmes d'accroche et d'assemblage, phases PRO et EXE.
    \item Gestion de sessions d'impression 3D grande échelle multimatériaux (béton, plâtre, terre, géopolymère).
    \item Formation à la fabrication robotique.
  \end{itemize}

  \item[2016 -- 2019] \textbf{BIM management}, freelance
  \begin{itemize}
    \item Accompagnement et implémentation de processus BIM : chartes BIM pour la MOE, cahiers des charges BIM pour la MOA, suivi des livrables.
    \item Vérification de maquettes numériques IFC (lycées, logements, bureaux), toutes phases de l'APS au DOE, marchés publics et privés.
    \item Production de maquettes interactives VR.
  \end{itemize}

  \item[2015 -- 2016] \textbf{BIM manager}, Astus Construction
  \begin{itemize}
    \item Rédaction de chartes BIM pour la MOE et de cahiers des charges BIM pour la MOA.
  \end{itemize}
\end{description}

% ------------------------------------------------------------------
% ------------------------------------------------------------------
\section*{Diplômes}

\begin{description}
  \item[2022 -- 2025] \textbf{École nationale des ponts et chaussées}, diplôme national de doctorat
  \begin{itemize}
    \item Thèse : \emph{Étude des procédés d'impression 3D de bétons fibrés pour la construction. Formulation, renforcement et application}. Préparée au laboratoire Navier, soutenue le 12 septembre 2025.
  \end{itemize}

  \item[2013 -- 2018] \textbf{INSA Lyon}, diplôme d'ingénieur
  \begin{itemize}
    \item Génie civil et urbanisme, spécialité bâtiment, structure et énergie.
  \end{itemize}

  \item[2011 -- 2016] \textbf{ENSA Lyon}, diplôme d'État d'architecte
  \begin{itemize}
    \item Spécialité architecture numérique et formes complexes.
    \item Sélectionné au prix de la jeune architecture.
  \end{itemize}
\end{description}

% ------------------------------------------------------------------


\end{document}
